\subsection{逆步表示}

设 \(E^{*}\) 是 E 的对偶。由于 W 通过 \(\sigma\) 作用于 E，它通过结构的转移也作用于 \(E^{*}\) 。相应的表示

\[
\sigma^ {*}: \mathrm{W} \to \mathbf {G L} (\mathrm{E} ^ {*})
\]

称为 \(\sigma\) 的逆步表示。我们有

\[
\sigma^ {*} (w) = ^ {t} \sigma \left(w ^ {- 1}\right) \text {   for   all   } w \in W.
\]

若 \(x^{*} \in \mathrm{E}^{*}\) 且 \(w \in \mathrm{W}\) ，用 \(w(x^{*})\) 表示 \(x^{*}\) 在 \(\sigma^{*}(w)\) 下的像。

若 \(s \in S\) ，用 \(A_s\) 表示由 \(x^* \in E^*\) （其中 \(x^*(e_s) > 0\) ）构成的集合。设 \(C\) 为诸 \(A_s\) （\(s \in S\) ）的交。当 \(S\) 有限时，\(C\) 是 \(E^*\) 中的一个开单形锥（§1，第 6 号）。

定理 1（Tits）. 若 \(w \in W\) 且 \(C \cap w(C) \neq \emptyset\) ，则 \(w = 1\) 。

我们立即指出这个定理的若干推论：

推论 1. 群 W 简单传递地作用于诸 \(w(\mathbf{C})\) （\(w \in \mathbf{W}\) ）所成的集合上。

这是显然的。

推论 2. 表示 \(\sigma\) 与 \(\sigma^{*}\) 都是单射的。

事实上，若 \(\sigma^{*}(w) = 1\) ，则 \(w(\mathrm{C}) = \mathrm{C}\) ，于是由定理得 \(w = 1\) 。\(\sigma\) 的单射性由 \(\sigma^{*}\) 的单射性得出。

推论 3. 若 S 有限，则 \(\sigma(\mathrm{W})\) 是 \(\mathbf{GL}(\mathrm{E})\) （赋予其典范李群结构）的离散子群；类似地，\(\sigma^{*}(\mathrm{W})\) 是 \(\mathbf{GL}(\mathrm{E}^{*})\) 的离散子群。

设 \(x^{*} \in C\) 。由 \(g \in \mathbf{GL}(E^{*})\) （满足 \(g(x^{*}) \in C\) ）构成的集合 U 是 \(\mathbf{GL}(E^{*})\) 中单位元的邻域；由定理，

\[
\sigma^ {*} (\mathrm{W}) \cap \mathrm{U} = \{1 \};
\]

从而 \(\sigma^{*}(\mathbf{W})\) 是 \(\mathbf{GL}(\mathbf{E}^{*})\) 的离散子群。由结构的转移可知，\(\sigma(\mathbf{W})\) 在 \(\mathbf{GL}(\mathbf{E})\) 中是离散的。

定理 1 的证明。

若 \(w \in W\) ，用 \(l(w)\) 表示 \(w\) 关于 S 的长度（第 IV 章，§1，第 1 号）。

我们将证明以下断言，其中 \(n\) 是一个 \(\geqslant 0\) 的整数：

\((P_n)\) 设 \(w \in W\) 满足 \(l(w) = n\) ，并设 \(s \in S\) 。则：

或者 \(w(\mathbf{C})\subset \mathbf{A}_s\)

或者 \(w(\mathbf{C}) \subset s(\mathbf{A}_s)\) 且 \(l(sw) = l(w) - 1\) 。

\((Q_n)\) 设 \(w \in W\) 满足 \(l(w) = n\) ，并设 \(s, s' \in S\) ，\(s \neq s'\) 。设 \(W_{s,s'}\) 为 \(W\) 的由 \(s\) 与 \(s'\) 生成的子群。存在 \(u \in W_{s,s'}\) 使得

\[
w (C) \subset u \left(A _ {s} \cap A _ {s ^ {\prime}}\right) \quad \text { and } \quad l (w) = l (u) + l \left(u ^ {- 1} w\right).
\]

这些断言当 n = 0 时是平凡的。我们按照如下模式对 n 作归纳来证明它们

\[
\left(\left(P _ {n}\right) \text {   and   } \left(Q _ {n}\right)\right) \Longrightarrow \left(P _ {n + 1}\right) \quad \text { and } \quad \left(\left(P _ {n + 1}\right) \text {   and   } \left(Q _ {n}\right)\right) \Longrightarrow \left(Q _ {n + 1}\right).
\]

\(\left((\mathrm{P}_n)\text{ 且 } (\mathrm{Q}_n)\right) \Longrightarrow (\mathrm{P}_{n+1})\) 的证明。

设 \(w \in W\) 满足 \(l(w) = n + 1\) ，并设 \(s \in S\) 。我们可以把 \(w\) 写成 \(w = s'w'\) 的形式，其中 \(s' \in S\) 且 \(l(w') = n\) 。若 \(s' = s\) ，则把 \((P_n)\) 应用于 \(w'\) 表明 \(w'(C) \subset A_s\) ，从而 \(w(C) \subset s(A_s)\) 且 \(l(sw) = l(w') = l(w) - 1\) 。若 \(s' \neq s\) ，则把 \((Q_n)\) 应用于 \(w'\) 表明存在 \(u \in W_{s,s'}\) 使得

\[
w ^ {\prime} (\mathrm{C}) \subset u \left(\mathrm{A} _ {s} \cap \mathrm{A} _ {s ^ {\prime}}\right) \text { and } l \left(w ^ {\prime}\right) = l (u) + l \left(u ^ {- 1} w ^ {\prime}\right).
\]

我们有 \(w(\mathbf{C}) = s'w'(\mathbf{C}) \subset s'u(\mathbf{A}_s \cap \mathbf{A}_{s'})\) 。

引理 1. 设 \(s, s' \in S\) 满足 \(s \neq s'\) ，并设 \(v \in W_{s,s'}\) 。则 \(v(A_s \cap A_{s'})\) 含于 \(A_s\) 或 \(s(A_s)\) 之一，且在第二种情形有 \(l(sv) = l(v) - 1\) 。

其证明将在第 5 号中给出。

我们把引理应用于元素 \(v = s'u\) 。有两种可能：或者

\[
  s ^ {\prime} u \left(\mathrm{A} _ {s} \cap \mathrm{A} _ {s ^ {\prime}}\right) \subset \mathrm{A} _ {s} \quad \text { and } a fortiori w (\mathrm{C}) \subset \mathrm{A} _ {s},
\]

或者

\[
  s ^ {\prime} u \left(\mathrm{A} _ {s} \cap \mathrm{A} _ {s ^ {\prime}}\right) \subset s \left(\mathrm{A} _ {s}\right) \quad \text { and } \quad a fortiori w (\mathrm{C}) \subset s \left(\mathrm{A} _ {s}\right).
\]

此外，在第二种情形，\(l(ss'u) = l(s'u) - 1\) 。由此

\[
\begin{array}{c} l (s w) = l (s s ^ {\prime} w ^ {\prime}) = l (s s ^ {\prime} u. u ^ {- 1} w ^ {\prime}) \leqslant l (s s ^ {\prime} u) + l (u ^ {- 1} w ^ {\prime}) \\ = l (s ^ {\prime} u) + l (u ^ {- 1} w ^ {\prime}) - 1 \leqslant l (w) - 1, \end{array}
\]

并且我们知道这蕴含 \(l(sw) = l(w) - 1\) 。

证明 \(((\mathrm{P}_{n + 1})\) 与 \((\mathrm{Q}_n))\Longrightarrow (\mathrm{Q}_{n + 1}).\)

设 \(w \in W\) 满足 \(l(w) = n + 1\) ，并设 \(s, s' \in S\) ，\(s \neq s'\) 。若 \(w(C)\) 含于 \(A_s \cap A_{s'}\) ，则条件 \((Q_{n+1})\) 以 \(u = 1\) 得到满足。若不然，例如设 \(w(C)\) 不含于 \(A_s\) 。由 \((P_{n+1})\) ，有 \(w(C) \subset s(A_s)\) 且 \(l(sw) = n\) 。把 \((Q_n)\) 应用于 \(sw\) ，存在 \(v \in W_{s,s'}\) 使得

\[
s w (C) \subset v \left(A _ {s} \cap A _ {s ^ {\prime}}\right) \quad \text { and } \quad l (s w) = l (v) + l \left(v ^ {- 1} s w\right).
\]

则

\[
w (\mathrm{C}) \subset s v \left(\mathrm{A} _ {s} \cap \mathrm{A} _ {s ^ {\prime}}\right)
\]

且

\[
\begin{array}{c} l (w) = 1 + l (s w) = 1 + l (v) + l (v ^ {- 1} s w) \\ \geqslant l (s v) + l ((s v) ^ {- 1} w) \geqslant l (w), \end{array}
\]

因此上面的不等式必为等式。由此可知 \((\mathbf{Q}_{n + 1})\) 以 \(u = sv\) 得到满足。

定理的证明。

设 \(w \in W\) 满足 \(w \neq 1\) 。我们可以把 \(w\) 写成 \(sw'\) 的形式，其中 \(s \in S\) 且 \(l(w') = l(w) - 1\) 。把 \((\mathrm{P}_n)\) 应用于 \(w'\) 与 \(n = l(w')\) ，得 \(w'(C) \subset A_s\) ，因为情形 \(w'(C) \subset s(A_s)\) 由于 \(l(sw') = l(w) = l(w') + 1\) 而被排除。于是 \(w(C) = sw'(C) \subset s(A_s)\) ，而由于 \(A_s\) 与 \(s(A_s)\) 不相交，我们有 \(C \cap w(C) = \varnothing\) 。证毕。

